% GENERATED FILE. Edit canonical lessons, then run npm run generate:books.
% canonical-source-hash: fnv1a64:0b60fcf12d42ba93
% canonical-lessons: HI-C317-rahasya, HI-C317-mukhbir, HI-C317-andaza, HI-C317-caunkna, HI-C317-hairan

\chapter{The Truth Comes Out}
\label{ch:hi-the-truth-comes-out}
\hlchaptermodality{317}

\begin{chapteropening}
\textbf{By the end of this chapter:} \emph{I can say what turned out to be true.}

Complete the last lesson of chapter 317: I can say what turned out to be true.
\end{chapteropening}

\section[\texorpdfstring{\hi{रहस्य} (rahasya)}{रहस्य (rahasya)}]{\hi{रहस्य} (rahasya) — a mystery; a secret}

\label{lesson:HI-C317-rahasya}



\begin{quote}
Before the new one: say the Hindi for a staff, then the Hindi for a curfew.
\end{quote}

\subsection*{You'll want to know: \hi{रहस्य}}
\textbf{\hi{रहस्य}} — \emph{rahasya} — \textquotedblleft{}a mystery; a secret\textquotedblright{}.

\begin{quote}
\hi{चोरी} \hi{का} \hi{रहस्य} \hi{पाँच} \hi{साल} \hi{तक} \hi{नहीं} \hi{खुला।}
\end{quote}

\begin{quote}
\emph{corī kā rahasya pā̃c sāl tak nahīṃ khulā.}
\end{quote}

\begin{quote}
The mystery of the theft stayed unsolved for five years.
\end{quote}

Masculine.

\begin{cousinweb}
\textbf{\hi{रहस्य}} (\emph{rahasya}) is Sanskrit \emph{rahasya}, \textquotedblleft{}secret, private\textquotedblright{} (Monier-Williams), from \emph{rahas}, solitude. A mystery \textbf{\hi{खुलता} \hi{है}}, opens, when it is solved.
\end{cousinweb}

\begin{grammarlens}[title={What you've built}]
One more word for telling what happened.
\end{grammarlens}

\subsection*{Your turn}
\begin{itemize}
\raggedright
  \item \emph{Say it:} \emph{rahasya}
  \item \emph{Say it:} \emph{rahasya}, once more
  \item \emph{Say it:} \emph{corī kā rahasya pā̃c sāl tak nahīṃ khulā.} — The mystery of the theft stayed unsolved for five years.
  \item \emph{Recall:} say the Hindi for a staff, then the Hindi for a curfew, then say \emph{rahasya} again
\end{itemize}

\begin{tcolorbox}[breakable,colback=teal!4,colframe=teal!35!black,title={Before you move on}]
What does \hi{रहस्य} mean? (\textquotedblleft{}A mystery; a secret\textquotedblright{}.) Say it once more.
\end{tcolorbox}
\section[\texorpdfstring{\hi{मुखबिर} (mukhbir)}{मुखबिर (mukhbir)}]{\hi{मुखबिर} (mukhbir) — an informer}

\label{lesson:HI-C317-mukhbir}



\begin{quote}
Before the new one: say the Hindi for a curfew, then the Hindi for a mystery.
\end{quote}

\subsection*{You'll want to know: \hi{मुखबिर}}
\textbf{\hi{मुखबिर}} — \emph{mukhbir} — \textquotedblleft{}an informer\textquotedblright{}.

\begin{quote}
\hi{एक} \hi{मुखबिर} \hi{ने} \hi{पुलिस} \hi{को} \hi{चोर} \hi{का} \hi{पता} \hi{बताया।}
\end{quote}

\begin{quote}
\emph{ek mukhbir ne pulis ko cor kā patā batāyā.}
\end{quote}

\begin{quote}
An informer told the police where the thief was.
\end{quote}

Masculine or feminine, by the person meant.

\begin{cousinweb}
\textbf{\hi{मुखबिर}} (\emph{mukhbir}) is usually traced, through Persian, to Arabic \emph{mukhbir}, \textquotedblleft{}one who brings news\textquotedblright{}, the same family as \textbf{\hi{ख़बर}}, news: someone who passes news to the police.
\end{cousinweb}

\begin{grammarlens}[title={What you've built}]
One more word for telling what happened.
\end{grammarlens}

\subsection*{Your turn}
\begin{itemize}
\raggedright
  \item \emph{Say it:} \emph{mukhbir}
  \item \emph{Say it:} \emph{mukhbir}, once more
  \item \emph{Say it:} \emph{ek mukhbir ne pulis ko cor kā patā batāyā.} — An informer told the police where the thief was.
  \item \emph{Recall:} say the Hindi for a curfew, then the Hindi for a mystery, then say \emph{mukhbir} again
\end{itemize}

\begin{tcolorbox}[breakable,colback=teal!4,colframe=teal!35!black,title={Before you move on}]
What does \hi{मुखबिर} mean? (\textquotedblleft{}An informer\textquotedblright{}.) Say it once more.
\end{tcolorbox}
\section[\texorpdfstring{\hi{अंदाज़ा} (andāzā)}{अंदाज़ा (andāzā)}]{\hi{अंदाज़ा} (andāzā) — a guess; an estimate}

\label{lesson:HI-C317-andaza}



\begin{quote}
Before the new one: say the Hindi for a mystery, then the Hindi for an informer.
\end{quote}

\subsection*{You'll want to know: \hi{अंदाज़ा}}
\textbf{\hi{अंदाज़ा}} — \emph{andāzā} — \textquotedblleft{}a guess; an estimate\textquotedblright{}.

\begin{quote}
\hi{मेरा} \hi{अंदाज़ा} \hi{सही} \hi{निकला।}
\end{quote}

\begin{quote}
\emph{merā andāzā sahī niklā.}
\end{quote}

\begin{quote}
My guess turned out to be right.
\end{quote}

Masculine.

\begin{cousinweb}
\textbf{\hi{अंदाज़ा}} (\emph{andāzā}) is usually traced to Persian \emph{andāza}, \textquotedblleft{}a measure\textquotedblright{}: a guess is measuring by eye. \textbf{\hi{अंदाज़ा} \hi{लगाना}} is to guess.
\end{cousinweb}

\begin{grammarlens}[title={What you've built}]
One more word for telling what happened.
\end{grammarlens}

\subsection*{Your turn}
\begin{itemize}
\raggedright
  \item \emph{Say it:} \emph{andāzā}
  \item \emph{Say it:} \emph{andāzā}, once more
  \item \emph{Say it:} \emph{merā andāzā sahī niklā.} — My guess turned out to be right.
  \item \emph{Recall:} say the Hindi for a mystery, then the Hindi for an informer, then say \emph{andāzā} again
\end{itemize}

\begin{tcolorbox}[breakable,colback=teal!4,colframe=teal!35!black,title={Before you move on}]
What does \hi{अंदाज़ा} mean? (\textquotedblleft{}A guess; an estimate\textquotedblright{}.) Say it once more.
\end{tcolorbox}
\section[\texorpdfstring{\hi{चौंकना} (cauṅknā)}{चौंकना (cauṅknā)}]{\hi{चौंकना} (cauṅknā) — to be startled}

\label{lesson:HI-C317-caunkna}



\begin{quote}
Before the new one: say the Hindi for an informer, then the Hindi for a guess.
\end{quote}

\subsection*{You'll want to know: \hi{चौंकना}}
\textbf{\hi{चौंकना}} — \emph{cauṅknā} — \textquotedblleft{}to be startled\textquotedblright{}.

\begin{quote}
\hi{सच} \hi{सुनकर} \hi{सब} \hi{लोग} \hi{चौंक} \hi{गए।}
\end{quote}

\begin{quote}
\emph{sac sunkar sab log cauṅk gae.}
\end{quote}

\begin{quote}
Everyone was startled to hear the truth.
\end{quote}

Past (she …): \hi{वह} \hi{चौंक} \hi{गई}, \emph{vah cauṅk gaī}.

\begin{cousinweb}
\textbf{\hi{चौंकना}} (\emph{cauṅknā}) is listed by CDIAL beside \textbf{\hi{चमकना}}, to flash, under \emph{*camakka}, a sudden movement (CDIAL 4676): a start of surprise is a sudden jerk of the body. Hindi writes it with the anusvara.
\end{cousinweb}

\begin{grammarlens}[title={What you've built}]
One more word for telling what happened.
\end{grammarlens}

\subsection*{Your turn}
\begin{itemize}
\raggedright
  \item \emph{Say it:} \emph{cauṅknā}
  \item \emph{Say it:} \emph{cauṅknā}, once more
  \item \emph{Say it:} \emph{sac sunkar sab log cauṅk gae.} — Everyone was startled to hear the truth.
  \item \emph{Recall:} say the Hindi for an informer, then the Hindi for a guess, then say \emph{cauṅknā} again
\end{itemize}

\begin{tcolorbox}[breakable,colback=teal!4,colframe=teal!35!black,title={Before you move on}]
What does \hi{चौंकना} mean? (\textquotedblleft{}To be startled\textquotedblright{}.) Say it once more.
\end{tcolorbox}
\section[\texorpdfstring{\hi{हैरान} (hairān)}{हैरान (hairān)}]{\hi{हैरान} (hairān) — astonished}

\label{lesson:HI-C317-hairan}



\begin{quote}
Before the new one: say the Hindi for a guess, then the Hindi for to be startled.
\end{quote}

\subsection*{You'll want to know: \hi{हैरान}}
\textbf{\hi{हैरान}} — \emph{hairān} — \textquotedblleft{}astonished\textquotedblright{}.

\begin{quote}
\hi{मुखबिर} \hi{का} \hi{नाम} \hi{सुनकर} \hi{दरोगा} \hi{हैरान} \hi{रह} \hi{गया।}
\end{quote}

\begin{quote}
\emph{mukhbir kā nām sunkar darogā hairān rah gayā.}
\end{quote}

\begin{quote}
The inspector was astonished to hear the informer's name.
\end{quote}

\begin{cousinweb}
\textbf{\hi{हैरान}} (\emph{hairān}) is usually traced, through Persian, to Arabic \emph{ḥayrān}, \textquotedblleft{}bewildered, at a loss\textquotedblright{}. \textbf{\hi{हैरान} \hi{रह} \hi{जाना}} is to be left speechless with surprise.
\end{cousinweb}

\begin{grammarlens}[title={What you've built}]
One more word for telling what happened.
\end{grammarlens}

\subsection*{Your turn}
\begin{itemize}
\raggedright
  \item \emph{Say it:} \emph{hairān}
  \item \emph{Say it:} \emph{hairān}, once more
  \item \emph{Say it:} \emph{mukhbir kā nām sunkar darogā hairān rah gayā.} — The inspector was astonished to hear the informer's name.
  \item \emph{Recall:} say the Hindi for a guess, then the Hindi for to be startled, then say \emph{hairān} again
\end{itemize}

\begin{tcolorbox}[breakable,colback=teal!4,colframe=teal!35!black,title={Before you move on}]
What does \hi{हैरान} mean? (\textquotedblleft{}Astonished\textquotedblright{}.) Say it once more.
\end{tcolorbox}
